\pagebreak
\section{Метод решения}
Сжимаются координатные блоки булева индекса ЛР7 --- списки документов термов (в булевом индексе координата --- это docID). Корпус --- 1\,000\,000 документов о туризме (ЛР1), индекс без сжатия занимает 1\,584 МиБ, из них списки --- 1\,388 МиБ (88\,\%): 364 млн записей по 4 байта. Остальные разделы (таблицы документов и термов, пул строк) не сжимались: задание --- про координатные блоки.

\subsection{Метод сжатия: разности + Variable Byte}
Список документов отсортирован по возрастанию, поэтому вместо docID хранятся разности соседних номеров: $g_1 = d_1$, $g_i = d_i - d_{i-1} \ge 1$. У частого терма разности маленькие (у терма, который есть в каждом документе, все равны 1), у редкого --- большие, до $10^6$. Каждая разность записывается кодом Variable Byte (VByte) из лекции: число делится на группы по 7 бит, начиная с младших; каждая группа занимает байт, старший бит байта $c$ --- признак конца числа, он равен 1 только в последнем байте.

\textbf{Побитовая схема одного числа} $v$ (биты $v_{31} \dots v_0$), байты в файле идут слева направо:
\begin{center}
\small
\begin{tabular}{|c|c|c|c|c|}
\hline
Байт 0 & Байт 1 & Байт 2 & Байт 3 & Байт 4 \\
\hline
$c\,v_6 \dots v_0$ & $c\,v_{13} \dots v_7$ & $c\,v_{20} \dots v_{14}$ & $c\,v_{27} \dots v_{21}$ & $c\,0\,0\,0\,v_{31} \dots v_{28}$ \\
\hline
\end{tabular}
\end{center}
Записывается столько байт, сколько нужно для старшей ненулевой группы: $v < 2^7$ --- 1 байт, $v < 2^{14}$ --- 2, $v < 2^{21}$ --- 3, $v < 2^{28}$ --- 4, иначе 5. Чтение: \texttt{v |= (b \& 0x7F) <{}< shift}, пока не встретится байт с $c = 1$.

Пример из лекции: список 824, 829, 215\,406 $\to$ разности 824, 5, 214\,577:
\begin{center}
\small
\begin{tabular}{|r|l|l|}
\hline
Разность & Группы по 7 бит (от младших) & Байты \\
\hline
824 & \texttt{0111000 0000110} & \texttt{0}\,\texttt{0111000} \texttt{1}\,\texttt{0000110} = \texttt{38 86} \\
5 & \texttt{0000101} & \texttt{1}\,\texttt{0000101} = \texttt{85} \\
214\,577 & \texttt{0110001 0001100 0001101} & \texttt{0}\,\texttt{0110001} \texttt{0}\,\texttt{0001100} \texttt{1}\,\texttt{0001101} = \texttt{31 0c 8d} \\
\hline
\end{tabular}
\end{center}
6 байт вместо 12. Байты разностей одного списка идут подряд без разделителей: конец числа виден по биту $c$, число чисел в списке --- $df$ из таблицы термов.

\textbf{Почему VByte.}
\begin{itemize}
    \item Код выровнен по байтам: распаковка --- маска, сдвиг и сложение на байт, без чтения отдельных битов через границы байтов, как у гамма-кода. Это самый простой и быстрый из кодов лекции, а поиск распаковывает список на каждый терм запроса.
    \item Выигрыш по размеру большой там, где сосредоточены данные. 73\,\% записей индекса --- у термов с $df > 10^4$, разности у них в основном меньше 128, т.\,е. 1 байт вместо 4. У редких термов разности до $10^6$ занимают 3 байта --- тоже меньше 4.
    \item По оценке в книге Маннинга~\cite{Kormen} (корпус RCV1) VByte уступает гамма-коду по размеру всего около 15\,\% (116 и 101 МБ при 400 МБ для 32-битных docID), а распаковывается быстрее.
\end{itemize}
Слабое место VByte --- минимум 8 бит на запись: у частых термов, где разность почти всегда 1, побитовые коды (гамма --- 1 бит на единицу) или битовые карты были бы в разы компактнее (см. выводы).

\subsection{Изменения формата индекса}
Формат BIDX версии 1 из ЛР7 расширен так, как было предусмотрено в ЛР7 полем \texttt{flags}:
\begin{itemize}
    \item \texttt{flags}, бит 1 (значение 2) --- \enquote{списки сжаты}. Бит 0 --- стемминг (ЛР4), флаги независимы: индекс со стеммингом и сжатием имеет \texttt{flags} = 3.
    \item В записи терма поле \texttt{first} (8 байт, смещение 8) при сжатии --- смещение списка в байтах от начала раздела списков, а не номер первой записи. Списки идут подряд в порядке термов, поэтому длина списка в байтах --- \texttt{first} следующего терма минус свой \texttt{first} (у последнего терма --- до начала пула строк, $\texttt{strings\_off} - \texttt{postings\_off}$).
    \item Раздел списков --- $\sum_t \mathrm{bytes}(t)$ байт кодов VByte вместо $P \times 4$; поле \texttt{P} заголовка по-прежнему --- число записей.
    \item Заголовок (кроме \texttt{flags} и \texttt{strings\_off}), таблица документов, поля \texttt{term} и \texttt{df} таблицы термов и пул строк не меняются. Терм по-прежнему ищется двоичным поиском по записям одного размера.
\end{itemize}
Читатель индекса теперь отказывается открывать файл с неизвестными флагами и проверяет, что смещения \texttt{first} не убывают и не выходят за раздел списков.

\subsection{Встраивание в поисковую систему}
\begin{itemize}
    \item \texttt{lab11/src/vbyte.h} --- кодек: длина кода, упаковка и распаковка числа и списка с проверками (конец списка, число длиннее 5 байт или больше $2^{32} - 1$, нулевая разность, переполнение docID).
    \item Индексатор ЛР7 получил ключ \texttt{-c}. Списки по-прежнему копятся в памяти массивами docID и сжимаются при записи файла. Таблица термов в файле идёт перед списками, поэтому сначала считаются длины кодов всех списков (\texttt{vb::list\_size}) и смещения \texttt{first}, затем списки пишутся кодами VByte прямо в буфер записи.
    \item Читатель индекса (\texttt{lab7/src/index\_reader.h}) получил метод \texttt{read\_list}: для сжатого индекса распаковывает список в массив docID, для несжатого --- копирует числа. Через него читают списки поиск ЛР8 и \texttt{lookup} ЛР7. Пересечение, объединение и разность в поиске не менялись: они работают с распакованными массивами, как раньше. Повреждённый список поиск считает пустым и сообщает об этом в stderr.
\end{itemize}

\subsection{Методика замеров}
Сравниваются два индекса одного корпуса: \texttt{lab7/data/index.bin} (без сжатия) и \texttt{lab11/data/index-vb.bin} (VByte).
\begin{itemize}
    \item \textbf{Проход по координатным блокам} --- утилита \texttt{build/listbench}: загружает индекс в память и для группы термов читает каждый список через \texttt{read\_list} в массив (как поиск) и суммирует docID. Время --- минимум из 10 проходов (2 запуска по 5). Группы по $df$: \emph{редкие} --- $df \le 100$ (меньше 0,01\,\% документов), \emph{средние} --- $100 < df \le 10^4$, \emph{частые} --- $df > 10^4$ (больше 1\,\% документов), и все термы.
    \item \textbf{Поиск} --- 24 запроса ЛР8 (\texttt{lab8/data/queries.txt}), утилита \texttt{search} с повтором каждого запроса 20 раз, минимум из 3 запусков. Отдельно --- время загрузки индекса в память.
    \item \textbf{Индексация} --- \texttt{indexer} и \texttt{indexer -c} на всём корпусе, по одному запуску (каждый около 9 минут), время по этапам.
\end{itemize}
